El estudio busca anticipar el estado con el que cada contratación queda registrada a partir del texto de su PBC. Para ello se utilizan únicamente procesos que cuentan con un documento aprovechable y con uno de los estados finales considerados. Los estados se agrupan en dos resultados:

\begin{itemize}
  \item Se considera exitoso el proceso registrado como \texttt{complete}.
  \item Se considera no exitoso el proceso registrado como \texttt{unsuccessful}, \texttt{cancelled} o \texttt{canceled}.
\end{itemize}

Los términos \texttt{cancelled} y \texttt{canceled} expresan el mismo estado, ``cancelado'', con las grafías británica y estadounidense del inglés; ambos aparecen en los datos de origen y por eso se normalizan dentro del mismo grupo. Reducir los estados a estas dos alternativas convierte el problema en una clasificación binaria: para cada PBC, el modelo estima la probabilidad de que el proceso pertenezca al primer grupo.

Los registros con otros estados se excluyen del entrenamiento y de la evaluación. Esta agrupación es una decisión operativa basada en la fuente disponible. Un proceso \texttt{complete} no implica necesariamente máximo valor público, y una cancelación puede responder a causas legítimas; por tanto, la etiqueta no constituye una evaluación integral de éxito, eficiencia o integridad.
